
%\lee{time permitting, things interesting to evaluate: (1) how personalized caption generation diffs per different domains--eg, CS vs. Bio papers; (2) how the types of target or profile figures affect the generation quality--eg line/bar/pie charts vs. other general illustration; (3) how the choice of profile figure affects the generation quality when there are multiple profile figures, or when target/profile figures are swapped, what happens?} \ed{I think we can calculate the first two?}\sam{calculated, but seems no big difference among them...}
%\sam{changed from "four" to "several", included the footnote underneath}

\paragraph{Experiment Setups.}\label{para:Experiment-Setups}
%\paragraph{Models.}
We evaluated four LLMs
%\footnote{We also evaluated the smaller Qwen-2.5-VL-7B-Instruct model, but excluded it from our main analysis due to a significantly higher rate of generation failures (n=402) compared to the other LLMs.} 
on personalized caption generation using \dataset:\footnote{Qwen-2.5-VL-7B-Instruct was excluded from our main analysis due to significantly higher failure rate (2.2\%, 269/12,259) compared to other models. See~\autoref{app:output-cleaning}.}
%\kenneth{TODO: In Appendix: (1) Mention Qwen's detailed error rate *as compared to other LLMs*, (2) provide examples of the invalid output, and (3) Still report the LLM performance numbers (do we have this?).}\sam{the Qwen one we only have the basic one with no profile, 1-profile, all-profile, but no "Other", also added the table there.}
{\em (i)} GPT-4o~\cite{hurst2024gpt}, 
{\em (ii)} Llama 4 Scout~\cite{meta2025llama4}, 
{\em (iii)} Gemini 2.5 Flash Preview~\cite{gemini25flash}, and 
{\em (iv)} GPT-4.1 Mini~\cite{openai2024gpt41mini}.
%\kenneth{TODO: Mention Qwen in the footnote ONLY.}
%\kenneth{TODO: Add citations to all.}
The first three are larger models, while the last one is smaller.
%We used OpenRouter (openrouter.ai), a unified API providing access to multiple AI models via a single endpoint, to prompt these models in a few-shot manner. 
We used OpenAI's API for GPT-4o and OpenRouter (openrouter.ai) for the others. 
We focused on LLMs because large-scale human evaluations from the \textsc{SciCap} Challenge showed a clear performance gap between model classes~\cite{hsu2025large}: only LLMs like GPT-4V consistently generate captions matching or exceeding those by human authors, while smaller or specialized models such as PEGASUS~\cite{zhang2020pegasus} and UniChart~\cite{masry2023unichart} perform poorly.
%\kenneth{TODO: Cite PEGASUS and UniChart}\sam{added}




%\sam{ADD-START}
%Recent large-scale human evaluations from the \textsc{SciCap} Challenge reveal a clear performance gap between model classes~\cite{hsu2025large}. These studies, judged by professional editors and PhD students, show that only state-of-the-art LLMs (\eg, GPT-4V) can generate captions that match or exceed the quality of those written by human authors. In contrast, smaller or more specialized models like PEGASUS and UniChart have performed poorly in such evaluations. This finding motivates our focus on the unique capabilities of modern LLMs for this task.
%\sam{ADD-END}





%\paragraph{Profile Setups.}
%All models were prompted in a zero-shot setting.
%Each model generated captions based on the target figure's image and figure-mentioning paragraphs, while also considering the profile figures' source-caption pairs for personalization.
%\tong{Might I ask why we focus on the three settings? Might it be helpful to cite some papers using similar settings?} 
%Building on findings that the amount of profile information affects performance~\cite{tan2024democratizing}, we tested three settings for caption generation for each model:
Building on prior work showing that more profile information improves performance~\cite{tan2024democratizing}, we tested four caption generation settings with varying amounts and sources of profile input:
\textbf{(1) No Profile:} The model generated captions using only the target figure's image and figure-mentioning paragraphs.
\textbf{(2) One Profile:} The model used the same source as in (1) but additionally used \textit{one} randomly selected profile figure from the \textit{same} paper
for personalization.
\textbf{(3) All Profile:} The model used the same source as in (1) but additionally used \textit{all} profile figures from the \textit{same} paper for personalization.
%\kenneth{My bad. Let's use ``Other Profile'' here. It's more direct.}
\textbf{(4) Other Profile:} 
The model used the same source as in (1) but additionally used
%Same as (1) plus 
one or three randomly selected profile figure(s) from random \textit{other} papers. 
The setup (4) tests whether performance gains come from paper-specific context 
%(personalization) 
or generic in-domain examples.
See~\autoref{app:prompts} for the full prompt.
%\kenneth{TODO: Add Unrelated profile prompts to the Appnedix.} \sam{indeed the prompt is the same as "prompt with profile", the process is just different when do the sampling of the papers and the profile figures}
%\footnote{We avoid naming our approach as ``zero-shot'' or ``few-shot'' because profile figures and captions serve as implicit examples.}
%\sam{ADD-START}
%To isolate the effect of personalization, we established a control condition using profile figures from randomly selected papers. This experiment tests whether performance gains stem from paper-specific context (personalization) or from generic, in-domain examples. By comparing results from same-paper versus cross-paper profiles, we can measure the impact of document-level style.
%\sam{ADD-END}

%\paragraph{Output \& Evaluation.}
%Following generation, 
We cleaned the output by removing unnecessary reasoning steps or explanations. 
%that were not part of the actual caption.
We also removed cases (56 out of 12,259) where models failed to generate valid output. 
%(\eg, ``I'm sorry, I can't assist with that request.'').
See \autoref{app:output-cleaning} and %shows data cleaning procedures.
\autoref{app:text-evaluation} for details.
%shows the evaluation packages used.

% \begin{table}
% \footnotesize
% \setlength{\tabcolsep}{2pt}
% \centering
% \begin{tabular}{@{}l@{\hspace{1pt}}l@{\hspace{1pt}}cccc@{\hspace{2pt}}ccc@{}}
% \toprule
% \multirow{2}{*}{\textbf{LLM}} & \multirow{2}{*}{\textbf{\begin{tabular}[c]{@{}c@{}}Profile\\ Used\end{tabular}}} & \multicolumn{4}{c}{\textbf{BLEU}} & \multicolumn{3}{c}{\textbf{ROUGE}} \\ 
% \cmidrule(lr){3-6} \cmidrule(l){7-9}
%  &  & \textbf{B-1} & \textbf{B-2} & \textbf{B-3} & \textbf{B-4} & \textbf{R-1} & \textbf{R-2} & \textbf{R-L} \\ 
% \midrule
% \multirow{3}{*}{\textbf{\begin{tabular}[c]{@{}l@{}}GPT-\\ 4o\end{tabular}}} & \textbf{No} & .219 & .133 & .091  & .063 & .321 & .127 & .248 \\ \cline{2-9} 
%  & \textbf{One} & .279 & .186 & .137 & .103 & .384 & .178 & .313 \\ \cline{2-9} 
%  & \textbf{All} & .292 & .200 & .150 & .115 & .397 & .194 & .328 \\ \hline
% \multirow{3}{*}{\textbf{\begin{tabular}[c]{@{}l@{}}llama \\ 4\\ scout\end{tabular}}} & \textbf{No} & .254 & .178 & .138 & .112 & .357 & .182 & .293 \\ \cline{2-9} 
%  & \textbf{One} & .372 & .293 & .246 & .211 & .481 & .300 & .423 \\ \cline{2-9} 
%  & \textbf{All} & .396 & .318 & .270 & .235 & .503 & .324 & .447 \\ \hline
% \multirow{3}{*}{\textbf{\begin{tabular}[c]{@{}l@{}}Gemini\\ 2.5\\ Flash\end{tabular}}} & \textbf{No} & .305 & .230 & .188 & .160 & .417 & .237 & .361 \\ \cline{2-9} 
%  & \textbf{One} & .370 & .291 & .244 & .209 & .482 & .301 & .426 \\ \cline{2-9} 
%  & \textbf{All} & .395 & .317 & .270 & .234 & .504 & .328 & .449 \\ \hline
% \multirow{3}{*}{\textbf{\begin{tabular}[c]{@{}l@{}}GPT-\\ 4o\\ mini\end{tabular}}} & \textbf{No} & .209 & .124 & .081 & .054 & .305 & .117 & .225 \\ \cline{2-9} 
%  & \textbf{One} & .286 & .202 & .155 & .121 & .398 & .207 & .326 \\ \cline{2-9} 
%  & \textbf{All} & .300 & .218 & .171 & .137 & .412 & .225 & .342 \\ \hline
% \end{tabular}
% \caption{Performance comparison of LLMs with different profile settings.\kenneth{TODO Ed and Sam: Figure out a way to make this Table fit in 1-column. We don't have space for 2-column.}\sam{a bit tight, but see if works}\kenneth{There must be ways we can fit this within 1-column. Maybe ask CY?}}
% \label{tab:main-result}
% \end{table}


\begin{table}[t]
\centering
\footnotesize
%\scriptsize
\renewcommand{\arraystretch}{1.215}
\setlength{\tabcolsep}{2.0pt}
\definecolor{vlightgreen}{rgb}{0.95, 1.0, 0.95}
\definecolor{lightgreen}{rgb}{0.85, 0.95, 0.85}
\definecolor{darkgreen}{rgb}{0.7, 0.9, 0.7}

\definecolor{lightred}{rgb}{1.0, 0.9, 0.9}
\definecolor{darkred}{rgb}{0.9, 0.7, 0.7}



% Please add the following required packages to your document preamble:
% \usepackage{booktabs}
% \usepackage{multirow}
% \usepackage[table,xcdraw]{xcolor}
% Beamer presentation requires \usepackage{colortbl} instead of \usepackage[table,xcdraw]{xcolor}



% Please add the following required packages to your document preamble:
% \usepackage{booktabs}
% \usepackage{multirow}
% \usepackage[table,xcdraw]{xcolor}
% Beamer presentation requires \usepackage{colortbl} instead of \usepackage[table,xcdraw]{xcolor}

% Please add the following required packages to your document preamble:
% \usepackage{booktabs}
% \usepackage{multirow}
% \usepackage[table,xcdraw]{xcolor}
% Beamer presentation requires \usepackage{colortbl} instead of \usepackage[table,xcdraw]{xcolor}


% Please add the following required packages to your document preamble:
% \usepackage{booktabs}
% \usepackage{multirow}
% \usepackage[table,xcdraw]{xcolor}
% Beamer presentation requires \usepackage{colortbl} instead of \usepackage[table,xcdraw]{xcolor}


% Please add the following required packages to your document preamble:
% \usepackage{booktabs}
% \usepackage{multirow}
% \usepackage[table,xcdraw]{xcolor}
% Beamer presentation requires \usepackage{colortbl} instead of \usepackage[table,xcdraw]{xcolor}


\begin{tabular}{@{}lccccccccc@{}}
\toprule
 &
\multicolumn{2}{c}{\textbf{Profile}} &
  \multicolumn{4}{c}{\textbf{BLEU}} &
  \multicolumn{3}{c}{\textbf{ROUGE}} \\ \cmidrule(l){2-10} 
\multirow{-2}{*}{\textbf{LLM}} &
  \textbf{\begin{tabular}[c]{@{}c@{}}Same\\ Paper\end{tabular}} &
  \textbf{\begin{tabular}[c]{@{}c@{}}\#\\ $\le 3$\end{tabular}} &
  \textbf{B-1} &
  \textbf{B-2} &
  \textbf{B-3} &
  \textbf{B-4} &
  \textbf{R-1} &
  \textbf{R-2} &
  \textbf{R-L} \\ \midrule
 &
  N/A &
  0 &
  .219 &
  .133 &
  .091 &
  .063 &
  .321 &
  .127 &
  .248 \\ \cmidrule(l){2-10} 
 &
   &
  \cellcolor{darkred}1 &
  \cellcolor{darkred}.189 &
  \cellcolor{darkred}.108 &
  \cellcolor{darkred}.070 &
  \cellcolor{darkred}.046 &
  \cellcolor{darkred}.289 &
  \cellcolor{darkred}.101 &
  \cellcolor{darkred}.218 \\
 &
  \multirow{-2}{*}{N} &
  \cellcolor{lightred}3 &
  \cellcolor{lightred}.206 &
  \cellcolor{lightred}.118 &
  \cellcolor{lightred}.076 &
  \cellcolor{lightred}.051 &
  \cellcolor{lightred}.300 &
  \cellcolor{lightred}.107 &
  \cellcolor{lightred}.223 \\ \cmidrule(l){2-10} 
 &
   &
  \cellcolor{lightgreen}1 &
  \cellcolor{lightgreen}.279 &
  \cellcolor{lightgreen}.186 &
  \cellcolor{lightgreen}.137 &
  \cellcolor{lightgreen}.103 &
  \cellcolor{lightgreen}.384 &
  \cellcolor{lightgreen}.178 &
  \cellcolor{lightgreen}.313 \\
\multirow{-5}{*}{\begin{tabular}[c]{@{}l@{}}GPT-\\ 4o\end{tabular}} &
  \multirow{-2}{*}{Y} &
    \cellcolor{darkgreen} All &
  \cellcolor{darkgreen}.292 &
  \cellcolor{darkgreen}.200 &
  \cellcolor{darkgreen}.150 &
  \cellcolor{darkgreen}.115 &
  \cellcolor{darkgreen}.397 &
  \cellcolor{darkgreen}.194 &
  \cellcolor{darkgreen}.328 \\ \midrule
 &
  N/A &
  0 &
  .254 &
  .178 &
  .138 &
  .112 &
  .357 &
  .182 &
  .293 \\ \cmidrule(l){2-10} 
 &
   &
  \cellcolor{darkred}1 &
  \cellcolor{darkred}.223 &
  \cellcolor{darkred}.167 &
  \cellcolor{darkred}.132 &
  \cellcolor{darkred}.107 &
  \cellcolor{darkred}.348 &
  \cellcolor{darkred}.180 &
  \cellcolor{darkred}.295 \\
 &
  \multirow{-2}{*}{N} &
  \cellcolor{vlightgreen}3 &
  \cellcolor{vlightgreen}.251 &
  \cellcolor{vlightgreen}.187 &
  \cellcolor{vlightgreen}.152 &
  \cellcolor{vlightgreen}.129 &
  \cellcolor{vlightgreen}.368 &
  \cellcolor{vlightgreen}.202 &
  \cellcolor{vlightgreen}.317 \\ \cmidrule(l){2-10} 
 &
   &
  \cellcolor{lightgreen}1 &
  \cellcolor{lightgreen}.372 &
  \cellcolor{lightgreen}.293 &
  \cellcolor{lightgreen}.246 &
  \cellcolor{lightgreen}.211 &
  \cellcolor{lightgreen}.481 &
  \cellcolor{lightgreen}.300 &
  \cellcolor{lightgreen}.423 \\
\multirow{-5}{*}{\begin{tabular}[c]{@{}l@{}}Llama-\\ 4 \\ Scout\end{tabular}} &
  \multirow{-2}{*}{Y} &
  \cellcolor{darkgreen}All &
  \cellcolor{darkgreen}.396 &
  \cellcolor{darkgreen}.318 &
  \cellcolor{darkgreen}.270 &
  \cellcolor{darkgreen}.235 &
  \cellcolor{darkgreen}.503 &
  \cellcolor{darkgreen}.324 &
  \cellcolor{darkgreen}.447 \\ \midrule
 &
  N/A &
  0 &
  .305 &
  .230 &
  .188 &
  .160 &
  .417 &
  .237 &
  .361 \\ \cmidrule(l){2-10} 
 &
   &
  \cellcolor{darkred}1 &
  \cellcolor{darkred}.268 &
  \cellcolor{darkred}.198 &
  \cellcolor{darkred}.159 &
  \cellcolor{darkred}.132 &
  \cellcolor{darkred}.388 &
  \cellcolor{darkred}.213 &
  \cellcolor{darkred}.335 \\
 &
  \multirow{-2}{*}{N} &
  \cellcolor{lightred}3 &
  \cellcolor{lightred}.279 &
  \cellcolor{lightred}.205 &
  \cellcolor{lightred}.164 &
  \cellcolor{lightred}.137 &
  \cellcolor{lightred}.398 &
  \cellcolor{lightred}.217 &
  \cellcolor{lightred}.343 \\ \cmidrule(l){2-10} 
 &
   &
  \cellcolor{lightgreen}1 &
  \cellcolor{lightgreen}.370 &
  \cellcolor{lightgreen}.291 &
  \cellcolor{lightgreen}.244 &
  \cellcolor{lightgreen}.209 &
  \cellcolor{lightgreen}.482 &
  \cellcolor{lightgreen}.301 &
  \cellcolor{lightgreen}.426 \\
\multirow{-5}{*}{\begin{tabular}[c]{@{}l@{}}Gemini-\\ 2.5\\ Flash \\ Preview\end{tabular}} &
  \multirow{-2}{*}{Y} &
  \cellcolor{darkgreen}All &
  \cellcolor{darkgreen}.395 &
  \cellcolor{darkgreen}.317 &
  \cellcolor{darkgreen}.270 &
  \cellcolor{darkgreen}.234 &
  \cellcolor{darkgreen}.504 &
  \cellcolor{darkgreen}.328 &
  \cellcolor{darkgreen}.449 \\ \midrule
 &
  N/A &
  0 &
  .209 &
  .124 &
  .081 &
  .054 &
  .305 &
  .117 &
  .225 \\ \cmidrule(l){2-10} 
 &
   &
  \cellcolor{vlightgreen}1 &
  \cellcolor{vlightgreen}.216 &
  \cellcolor{vlightgreen}.133 &
  \cellcolor{vlightgreen}.091 &
  \cellcolor{vlightgreen}.063 &
  \cellcolor{vlightgreen}.325 &
  \cellcolor{vlightgreen}.132 &
  \cellcolor{vlightgreen}.250 \\
 &
  \multirow{-2}{*}{N} &
  \cellcolor{vlightgreen}3 &
  \cellcolor{vlightgreen}.220 &
  \cellcolor{vlightgreen}.136 &
  \cellcolor{vlightgreen}.093 &
  \cellcolor{vlightgreen}.065 &
  \cellcolor{vlightgreen}.328 &
  \cellcolor{vlightgreen}.134 &
  \cellcolor{vlightgreen}.252 \\ \cmidrule(l){2-10} 
 &
   &
  \cellcolor{lightgreen}1 &
  \cellcolor{lightgreen}.286 &
  \cellcolor{lightgreen}.202 &
  \cellcolor{lightgreen}.155 &
  \cellcolor{lightgreen}.121 &
  \cellcolor{lightgreen}.398 &
  \cellcolor{lightgreen}.207 &
  \cellcolor{lightgreen}.326 \\
\multirow{-5}{*}{\begin{tabular}[c]{@{}l@{}}GPT-\\ 4.1 \\ Mini\end{tabular}} &
  \multirow{-2}{*}{Y} &
  \cellcolor{darkgreen}All &
  \cellcolor{darkgreen}.300 &
  \cellcolor{darkgreen}.218 &
  \cellcolor{darkgreen}.171 &
  \cellcolor{darkgreen}.137 &
  \cellcolor{darkgreen}.412 &
  \cellcolor{darkgreen}.225 &
  \cellcolor{darkgreen}.342 \\ \bottomrule
\end{tabular}


%\caption{Performance of LLMs on caption generation using different profile settings. The results show that using profile figures (especially all available ones) consistently improves performance across all models.}
\vspace{-.5pc}
\caption{Performance of LLMs on caption generation across profile settings. The highest scores are achieved by using \textbf{all} available profile(s) from the \textbf{same} paper.}
\vspace{-1pc}
\label{tab:llm_comparison}
\end{table}

%\kenneth{(1) Try to use \texttt{centering} tag and \texttt{footnotesize} tag to make it smaller-ish. (2) Change ``Paper'' to ``Source Paper'', sharing the same header of ``Profile'' with (4). (4) Change ``Profile (max.3)'' to  ``\#~Profile (max=3)'', sharing the same header of ``Profile'' with (3) (See the image I sent you.) (5) Replace ``Random'' with ``Unrelated'' or ``Unrel.''.}
%\sam{(1) Footnotesize is used already in previous version, I used a smaller size (scriptsize) to see if it works.}\kenneth{This feels too small to be honest. Let's use footnotesize. Try (1) replacing ``\#~Profile'' with just ``\#'', and (2) replace ``Source Paper'' with ``Same Paper'' and use "N/A", "Y", "N" as colum values.}

%\caption{Performance of LLMs on caption generation across different profile settings. The highest scores are achieved by using \textbf{all} available profile figures from the \textbf{same} paper.}


\paragraph{Using profile information (from the same paper) makes captions more similar to ground truth, especially with all profile figures.}\label{para:profile-result}
\autoref{tab:llm_comparison} shows the personalized caption generation results of four LLMs, evaluated using BLEU~\cite{papineni2002bleu} and ROUGE~\cite{lin2004rouge}.\footnote{We also explored BERTScore~\cite{zhang2019bertscore}, which correlated highly with BLEU and ROUGE. See \autoref{app:main-study-full}.}
%\kenneth{TODOs: (1) Add ref, (2) Add BERTScore results to Appendix.}\sam{Done!}
We used reference-based metrics to measure how closely the generated captions matched the original author-written captions, following a standard evaluation approach for personalized text generation used in well-known work like LongLaMP~\cite{kumar2024longlamp}.
The results show that incorporating profile information consistently improves caption quality across all four models.
Additionally, using all profile figures provides better results than using just one.
%\sam{The detailed distribution of the performance of BLEU-4 and ROUGE-2 across different models and profile configurations can be found in \autoref{app:main-study-full}.}
%See detailed  performance distributions across models and profile settings in \autoref{app:main-study-full}.
See \autoref{app:main-study-full} for details.
%Detailed performance distributions across models and profile configurations are in \autoref{app:main-study-full}.


% \renewcommand\thesubfigure{1\alph{subfigure}}


\begin{figure*}[t]
    \centering
    % Top row
    \begin{subfigure}{0.495\textwidth}
        \renewcommand\thesubfigure{a-1}
        \includegraphics[width=\textwidth]{figures/fig04_easyHard/01_easy_bleu4-v2.png}
        \vspace{-1.7pc}
        \caption{BLEU-4 (Context-Aligned Subset)}
        \label{fig:top_left}
    \end{subfigure}
    \hfill
    \begin{subfigure}{0.495\textwidth}
        \renewcommand\thesubfigure{b-1}
        \includegraphics[width=\textwidth]{figures/fig04_easyHard/03_hard_bleu4-v2.png}
        \vspace{-1.7pc}
        \caption{BLEU-4 (Context-Misaligned Subset)}
        \label{fig:top_right}
    \end{subfigure}
    
    % Bottom row
    \begin{subfigure}{0.495\textwidth}
         \renewcommand\thesubfigure{a-2}
        \includegraphics[width=\textwidth]{figures/fig04_easyHard/02_easy_rouge2-v4.png}
        \vspace{-1.7pc}
        \caption{ROUGE-2 (Context-Misaligned Subset)}
        \label{fig:bottom_left}
    \end{subfigure}
    \hfill
    \begin{subfigure}{0.495\textwidth}
        \renewcommand\thesubfigure{b-2}
        \includegraphics[width=\textwidth]{figures/fig04_easyHard/04_hard_rouge2-v4.png}
        \vspace{-1.7pc}
        \caption{ROUGE-2 (Context-Misaligned Subset)}
        \label{fig:bottom_right}
    \end{subfigure}
    %\caption{Comparison of the BLEU-4 and ROUGE-2 metrics on \dataset Context-Aligned and Context-Misaligned Subset across LLMs and profile configurations. Personalization is more effective when profile captions are highly similar to the target caption.
    %\kenneth{(1) The subfigure on the bottom row might not need labels, (2) ROUGE-2's figures can reduce heights, (3) can we also reduce the padding empty space on the left/right of each subfigure.}\sam{updated!}
    %}
    \vspace{-1.5pc}
    \caption{BLEU-4 and ROUGE-2 scores on \dataset's Context-Aligned and Context-Misaligned subsets, showing that personalization is most effective when profile caption are similar to the target (a-1, a-2).}
    \vspace{-1pc}
    \label{fig:all_figures}
    %\vspace{-3mm}  % Reduce space
\end{figure*}



%    \caption{BLEU-4 and ROUGE-2 scores on \dataset's Context-Aligned and Context-Misaligned subsets across LLMs and profile settings. Personalization is more effective when at least one profile caption closely matches the target caption (a-1, a-2).}

%\kenneth{The following is newly added. Please read.}
\paragraph{Using profiles from other papers often lowered BLEU and ROUGE scores, though there were exceptions.} 
\autoref{tab:llm_comparison} suggests that performance gains primarily came from paper-specific context, rather than generic in-domain examples. 
Profiles from other papers generally hurt performance, but some models, such as GPT-4.1 Mini, showed slight improvements. 
%In general, 
Furthermore, using more \textit{other} profiles tended to reduce the performance drop or, occasionally, provide minor gains.



\begin{table}
\centering
\footnotesize
\renewcommand{\arraystretch}{1.2}
\setlength{\tabcolsep}{2pt}
\definecolor{lightgreen}{rgb}{0.9, 1.0, 0.9}
\definecolor{darkgreen}{rgb}{0.7, 0.9, 0.7}
\begin{tabular}{lcccccccc}
\toprule
\multirow{2}{*}{\textbf{LLM}} & \multirow{2}{*}{\textbf{Same}} & \multicolumn{4}{c}{\textbf{BLEU}} & \multicolumn{3}{c}{\textbf{ROUGE}} \\ 
\cmidrule(lr){3-6} \cmidrule(lr){7-9}
 & \textbf{Type} & \textbf{B-1} & \textbf{B-2} & \textbf{B-3} & \textbf{B-4} & \textbf{R-1} & \textbf{R-2} & \textbf{R-L} \\ 
\midrule
\multirow{2}{*}{GPT-4o} 
 & No & .233 & .142 & .097 & .069 & .340 & .134 & .270 \\ 
 & \cellcolor{darkgreen}Yes & \cellcolor{darkgreen}.302 & \cellcolor{darkgreen}.208 & \cellcolor{darkgreen}.157 & \cellcolor{darkgreen}.121 & \cellcolor{darkgreen}.406 & \cellcolor{darkgreen}.201 & \cellcolor{darkgreen}.336 \\ 
\midrule
\multirow{2}{*}{Llama-4 Scout}
 & No & .325 & .245 & .199 & .167 & .436 & .250 & .377 \\ 
 & \cellcolor{darkgreen}Yes & \cellcolor{darkgreen}.396 & \cellcolor{darkgreen}.317 & \cellcolor{darkgreen}.269 & \cellcolor{darkgreen}.233 & \cellcolor{darkgreen}.504 & \cellcolor{darkgreen}.325 & \cellcolor{darkgreen}.447 \\ 
\midrule
\multirow{2}{*}{\begin{tabular}[c]{@{}l@{}}Gemini-2.5\\Flash Preview\end{tabular}} 
 & No & .326 & .246 & .201 & .169 & .440 & .254 & .384 \\ 
 & \cellcolor{darkgreen}Yes & \cellcolor{darkgreen}.393 & \cellcolor{darkgreen}.314 & \cellcolor{darkgreen}.266 & \cellcolor{darkgreen}.229 & \cellcolor{darkgreen}.503 & \cellcolor{darkgreen}.325 & \cellcolor{darkgreen}.447 \\ 
\midrule
\multirow{2}{*}{GPT-4.1 Mini}  
 & No & .237 & .154 & .109 & .079 & .349 & .155 & .276 \\ 
 & \cellcolor{darkgreen}Yes & \cellcolor{darkgreen}.311 & \cellcolor{darkgreen}.227 & \cellcolor{darkgreen}.178 & \cellcolor{darkgreen}.142 & \cellcolor{darkgreen}.422 & \cellcolor{darkgreen}.234 & \cellcolor{darkgreen}.351 \\ 
\bottomrule
\end{tabular}
%\caption{Performance comparison of LLMs with 1 profile settings. Comparing the data with same figure type (n=8083) and different figure type (n=4120). When profile figures shared the same type as the target figure, personalization works better.}
\vspace{-.5pc}
\caption{LLM performance on figures with one profile figure. Personalization is more effective when the single profile figure shares the same type as the target.}
\vspace{-1pc}
\label{tab:main-result-sameType}
\end{table}




\paragraph{When profile figures shared the same type as the target figure, personalization works better.}
To examine how figure type affects personalization, we analyzed cases with a single profile figure, splitting them into two groups: those where the profile and target figure types matched (n=8,083) and those where they did not (n=4,120). \autoref{tab:main-result-sameType} shows that matching the figure type resulted in captions that were significantly closer to the gold caption.

% The \textsc{SciCap} Challenge dataset provided figure types (\eg, graph plot, scatterplot), allowing us to examine how figure types affect personalized caption generation. 
% We focused on cases where a target figure had only one profile figure, dividing them into two groups: those where the profile figure type matched the target figure (Yes) (n=8,083) and those where it did not (No) (n=4,120). 
% \autoref{tab:main-result-sameType} shows that captions were closer to the ground-truths when the profile figure shared the same type as the target figure.


\paragraph{Personalization is more effective when profile captions are highly similar to the target caption.}\label{para:context-alignment}

To test if personalization is more effective when profiles are similar to the target, we split our test set into two groups. We calculated the similarity (using BERTScore and ROUGE-L) between each target caption and its available profile captions. The top 25\% of examples with the most similar profiles formed our Context-Aligned set (n=2,513); the remainder formed the Context-Misaligned set. The results in \autoref{fig:all_figures} confirm our hypothesis. Performance gains from using profiles were substantially larger for the Context-Aligned group~(Figures~\ref{fig:top_left} and~\ref{fig:bottom_left}), while the impact was noticeably smaller for the Context-Misaligned set~((Figures~\ref{fig:top_right} and~\ref{fig:bottom_right}). 
(See \autoref{app:easy-hard}.)




% Inspired by the figure type study (\autoref{tab:main-result-sameType}), we investigated whether profile captions similar to the target caption improve personalization. 
% For each target figure in the test set, we calculated the semantic similarity using BERTScore~\cite{Zhang*2020-xq} and lexical similarity using ROUGE-L~\cite{Lin2004-ox} between profile captions and the target caption.
% (\autoref{app:easy-hard} shows detailed score distributions.)
% %\sam{(\autoref{fig:profile_target_distribtions})}\sam{just wonder should I add the figure directly from appendix to here? Or I should always refer to Appendix}.
% We then sorted the test set by the highest similarity score for each target figure and selected the top 25\% of data points with high similarity in both metrics, creating the Context-Aligned set with 2,513 target figures.
% The remaining 9,690 target figures in the test set formed the Context-Misaligned set.
% \autoref{fig:all_figures} shows that personalization is more effective when at least one profile caption is highly similar to the ground-truth caption (Figures~\ref{fig:top_left} and~\ref{fig:bottom_left}).
% While profiles still help in the Context-Misaligned set (Figures~\ref{fig:top_right} and~\ref{fig:bottom_right}), their impact is noticeably smaller. 
% The detailed results can be found in \autoref{app:easy-hard}.
\begin{table}[t]
\centering
\footnotesize
\renewcommand{\arraystretch}{1.2}
\setlength{\tabcolsep}{2pt}
\definecolor{lightgreen}{rgb}{0.9, 1.0, 0.9}
\definecolor{darkgreen}{rgb}{0.7, 0.9, 0.7}

% Please add the following required packages to your document preamble:
% \usepackage{booktabs}
% \usepackage{multirow}

\begin{tabular}{@{}clllllll@{}}
\toprule
\multirow{2}{*}{\begin{tabular}[c]{@{}c@{}}Profile\\ \#\end{tabular}} & \multicolumn{4}{c}{BLEU}                                                                              & \multicolumn{3}{c}{ROUGE}                                                   \\ \cmidrule(l){2-8} 
                                                                      & \multicolumn{1}{c}{B-1} & \multicolumn{1}{c}{B-2} & \multicolumn{1}{c}{B-3} & \multicolumn{1}{c}{B-4} & \multicolumn{1}{c}{R-1} & \multicolumn{1}{c}{R-2} & \multicolumn{1}{c}{R-L} \\ \midrule
0                                                                     & .212                    & .124                    & .082                     & .055                    & .302                    & .113                    & .223                    \\
1                                                                     & .289                    & .198                    & .145                    & .109                    & .390                    & .181                    & .312                    \\
2                                                                     & .319                    & .215                    & .159                    & .121                    & .411                    & .198                    & .331                    \\
3                                                                     & .332                    & .231                    & .175                    & .136                    & .424                    & .215                    & .345                    \\ \bottomrule
\end{tabular}

%\caption{Performance of LLMs on caption generation using different profile settings. The results show that using profile figures (especially all available ones) consistently improves performance across all models.}
\vspace{-.5pc}
%\caption{Performance of GPT-4o on caption generation with varying number of profile figures used.}
\caption{GPT-4o performance with varying numbers of profile figures. Scores improve as more profiles are added, with the largest gain from 0 to 1 profile.}
\vspace{-1pc}
\label{tab:4o_profile-number}
\end{table}






\paragraph{Adding more profile figures consistently improved performance, but the largest gain came from the first profile.}
We examined how the number of profiles affected caption quality using GPT-4o. 
A subset figure with exactly three profiles from the same paper was used (n=3,424). 
For the 1-Profile and 2-Profile settings, profiles were randomly sampled. 
\autoref{tab:4o_profile-number} showed clear diminishing returns as more profiles were added.


% \sam{ADD-END}



%\tong{Ablation Study seems to be Section 4.1. Might it be helpful to create some subsections for the above content?}
\subsection{Ablation Study\label{sec:ablation}}
% \paragraph{Caption is the most important thing; image next. Paragraph is the least important---even our INLG paper said so! (CITATION). \cite{huang-etal-2023-summaries}}
% \paragraph{Caption is the most important thing; image next. Paragraph is the least important---even our INLG paper said so! (CITATION). \cite{huang-etal-2023-summaries}}


\paragraph{Captions are the most important profile element, while images are more influential than paragraphs.}
To assess the importance of each profile element, we conducted an ablation study on the test set using the GPT-4o model with the One Profile setting.
We tested three conditions by removing one profile element at a time: 
{\em (i)} figure captions (No Caption), 
{\em (ii)} figure images (No Image), and 
{\em (iii)} figure-mentioning paragraphs (No Paragraph). 
\autoref{fig:ablation} shows the results. 
Removing captions had the most significant impact, as captions directly guide generation. 
Removing images also reduced performance more than removing paragraphs, highlighting the greater influence of visual information.
%\sam{The detailed result can be referred to \autoref{app:ablation-study-detail}.}
\autoref{app:ablation-study-detail} shows the detailed results.


%------------ dead kitten ----------
\begin{comment}

We conducted an ablation study on the test set (consisting of 12,270 target figures) to evaluate the impact of different profile components on caption generation performance using GPT-4o model. We tested three conditions by removing individual profile elements: (1) No Captions, (2) No Image, and (3) No Paragraph. Generated captions were compared against original target captions using BLEU and ROUGE scores. Results indicated that the 'No Captions' condition yielded the lowest scores, highlighting the critical role of profile captions in producing contextually accurate captions. The 'No Image' condition showed slightly lower performance compared to 'No Paragraph', suggesting images have a moderately greater impact than paragraphs. The detailed results are shown in Figure~\ref{fig:ablation}. These findings align with previous research emphasizing the importance of captions in generating coherent and contextually relevant text \cite{huang-etal-2023-summaries}.



\kenneth{TODO: Add a small table here}
\sam{added barchart above}\sam{not sure if still need table, maybe put in appendix?}

    
\end{comment}





\begin{figure}[t]
    \centering
    \includegraphics[width=\columnwidth]{figures/fig01/fig01_ablation_study-v5.png}
    %\caption{Result of ablation study on caption generation. Captions are the most important profile element, while images are more influential than paragraphs.
    %\kenneth{I like this design! (1) Enlarge font size. It's hard to read. (2) Can we pick a slightly darker color? These colors are fine, but a bit too light.}\sam{Updated with darker color and a bit larger font!}\kenneth{(1) What is the number on top of each bar? Maybe remove them. (2) Fonts need to be even bigger. (3) The label can be ``No Paragraph'', ``No Image''. ``No Caption'', (4) The bars can be slightly wider.}\sam{The number above is the average score for each setting. I think keeping them can make it more clear to know the actual difference? (Or create another table to show the full result of ablation study for clarity?)}\sam{P.S. I just make the chart flattened a bit, same setting as for EasyHard one}
    %}
    \vspace{-2pc}
    \caption{Ablation study on the impact of profile elements using GPT-4o. Results show a clear hierarchy of importance: caption > image > paragraph.}
    \vspace{-.5pc}
    \label{fig:ablation}
    %\vspace{-3mm}  % Reduce space
\end{figure}


%     \caption{Ablation study on the impact of profile elements. caption generation by removing one profile element at a time (caption, image, or figure-mentioning paragraph). Captions contribute the most, followed by images, then paragraphs.}

%\sam{added for camera-ready}
%\kenneth{NEWLY ADDED!}

%\subsection{Impact of Profile Quantity\label{sec:quantity}}
%\input{sections/4-2-selection-strategy} 

\subsection{Human Evaluation\label{sec:humanEval}}
% \paragraph{Caption is the most important thing; image next. Paragraph is the least important---even our INLG paper said so! (CITATION). \cite{huang-etal-2023-summaries}}
% \paragraph{Caption is the most important thing; image next. Paragraph is the least important---even our INLG paper said so! (CITATION). \cite{huang-etal-2023-summaries}}

% \begin{figure}[t]
    \centering
    \includegraphics[width=\columnwidth]{figures/fig01/fig01_ablation_study-v5.png}
    %\caption{Result of ablation study on caption generation. Captions are the most important profile element, while images are more influential than paragraphs.
    %\kenneth{I like this design! (1) Enlarge font size. It's hard to read. (2) Can we pick a slightly darker color? These colors are fine, but a bit too light.}\sam{Updated with darker color and a bit larger font!}\kenneth{(1) What is the number on top of each bar? Maybe remove them. (2) Fonts need to be even bigger. (3) The label can be ``No Paragraph'', ``No Image''. ``No Caption'', (4) The bars can be slightly wider.}\sam{The number above is the average score for each setting. I think keeping them can make it more clear to know the actual difference? (Or create another table to show the full result of ablation study for clarity?)}\sam{P.S. I just make the chart flattened a bit, same setting as for EasyHard one}
    %}
    \vspace{-2pc}
    \caption{Ablation study on the impact of profile elements using GPT-4o. Results show a clear hierarchy of importance: caption > image > paragraph.}
    \vspace{-.5pc}
    \label{fig:ablation}
    %\vspace{-3mm}  % Reduce space
\end{figure}


%     \caption{Ablation study on the impact of profile elements. caption generation by removing one profile element at a time (caption, image, or figure-mentioning paragraph). Captions contribute the most, followed by images, then paragraphs.}
%\kenneth{TODO Human Evaluation here: https://openreview.net/forum?id=ghxUYlolwr\#discussion}


%\input{table/table_human rank}
\begin{figure}[t]
    \centering
    \includegraphics[width=\columnwidth]{figures/fig_humanRankDist/humanRankDist-v11.png}
    \vspace{-1.7pc}
    \caption{Distribution of human preference rankings (lower is better). A Friedman's test confirmed a statistically significant difference between the model configurations ($\chi^2$ = 28.48, $p < 0.001$).}
    \vspace{-1pc}
    \label{fig:human-rank}
    %\vspace{-3mm}  % Reduce space
\end{figure}

%\kenneth{(1) I'd use capitalized P for No-Profile. I know no-profile technically is one word, BUT I'd like to keep things consistent. (2) Can we try a bar chart (each model has one bar)? Like the one we always did for human evaluation.}\sam{corrected!}
%\kenneth{(1) Img resolution is a bit low. Use a better quality image. (2) The Rank 1, 2, 3, 4 legend/label is too close to the bar. Hard to see.}\sam{updated}

%\kenneth{I'd use All-Profile, No-Profile instead of All-profile, No-profile. Use upper case P, as they are basically names of conditions.}

%\kenneth{TODO: Which model do we use in human eval????}

%\sam{ADD-START}
%To supplement automatic metrics, which do not always correlate with quality, we conducted a human evaluation.
We recruited 10 US-based NLP researchers (PhD students with publication or review experience) for our human evaluation.
%\kenneth{Be more specific, are they PhD students? or?}\sam{corrected!} 
The study involved ranking captions for 50 figures randomly selected from the arXiv \textit{cs.CL} domain. 
For each figure, we provided the figure's image, the paper's title, and abstract as context. 
Participants then ranked four corresponding captions based on how well each one helped them understand the figure. The four captions were:
%\kenneth{What exactly is the QUESTION we asked the participants? Don't just say QUALITY.} \sam{corrected! Also included the full instruction in the Appendix as requested in the EMNLP checklist}
{\em (i)} the original \textit{Gold} caption, and
{\em (ii)} \textit{No-Profile}, 
{\em (iii)} \textit{1-Profile}, and 
{\em (iv)} \textit{All-Profile} settings generated by our best-performing model (\textbf{Gemini}, based on automatic metrics). 
See \autoref{app:human-eval} for details.
%The complete experimental setup, which follows prior work~\cite{hsu2025large}, is detailed in \autoref{app:human-eval}.
%with the complete experimental setup and participant instructions detailed in \autoref{app:human-eval}.
%The ranking results are presented in 

\paragraph{1-Profile was the most preferred condition.} Our human evaluation result shows a clear preference for \textit{1-Profile} setting, which achieved the best average rank of 2.27. The other models followed in order: \textit{No-Profile} (2.48) and \textit{All-Profile} (2.54). Interestingly, the author-written \textit{Gold} captions were ranked last overall with an average rank of 2.71. Such outcome is reflected in the preference distribution (\autoref{fig:human-rank}): The \textit{1-Profile} configuration received the most first-place votes (30.6\%) and the fewest last-place votes (17.8\%). On contrast, the \textit{Gold} captions were ranked last most frequently (31.2\%).
%\kenneth{TODOs: As we do not have a table now, write a few sentences here to report the avg rank for each model here.}\sam{done!}
Post-hoc Friedman-Nemenyi tests showed that 1-Profile and No-Profile captions were significantly preferred over gold ($p < 0.001$ and $p = 0.03$, respectively).

\paragraph{Trade-offs between human-perceived quality and similarity to gold captions.} 
While the All-Profile setting generated captions with higher reference similarity (\autoref{tab:llm_comparison}), human judges significantly preferred captions from 1-Profile ($p = 0.006$). 
This suggests that optimizing for similarity may also reproducing flaws from the inconsistently quality's caption from arXiv reference~\cite{huang-etal-2023-summaries}, reducing the perceived quality.

%This may stem from the variable quality of the reference arXiv captions themselves~\cite{huang-etal-2023-summaries}. 
%Consequently, when a model too closely mimics the gold standard, it may inherit stylistic flaws, causing perceived quality to decline as similarity increases excessively.

%\sam{ADD-END}
%------------ dead kitten ----------
\begin{comment}


    %While the All-Profile setting generated captions with higher reference similarity, with higher similarity to the references (\autoref{tab:llm_comparison}), human judges significantly preferred captions from 1-Profile ($p = 0.006$). 

    % This likely reflects the inconsistent quality of arXiv captions~\cite{huang-etal-2023-summaries}.
Models that closely mimic these references risk reproducing their flaws, reducing perceived quality as similarity increases.
\end{comment}



 







%\sam{ADD-START}
%A semantic similarity analysis using BERTScore between gold caption and the generated captions confirmed these trends from additional profiles.\kenneth{What does this respond to?}\sam{This part is added in the rebuttal that we did BERTScore analysis and show correlation to the BLEU and ROUGE score}
%\sam{ADD-END}

% \sam{ADD-START}
% \paragraph{BERTScore results correlate well with our existing metrics.}
% We also report BERTScore in \autoref{tab:BERTScore_LLM} to evaluate semantic similarity between the generated and gold captions. The results correlate strongly with our BLEU and ROUGE metrics, confirming the same performance trends. The \textbf{All-Profile} setting achieves the highest score, followed closely by \textbf{1-Profile}. Notably, this analysis reveals a pattern of diminishing returns. The performance leap from No-Profile to 1-Profile is substantially larger than the subsequent gain from 1-Profile to All-Profile, indicating that while the first profile provides a significant semantic boost, the marginal benefit of additional profiles is less substantial.

































%----------dead kitten--------------
\begin{comment}


Inspired by the figure type study (\autoref{tab:main-result-sameType}), we investigated whether profile captions similar to the target caption improve personalization. 
For each target figure in the test set, we calculated the semantic similarity using BERTScore~\cite{Zhang*2020-xq} and lexical similarity using ROUGE-L~\cite{Lin2004-ox} between profile captions and the target caption.
(\autoref{app:easy-hard} shows detailed score distributions.)
%\sam{(\autoref{fig:profile_target_distribtions})}\sam{just wonder should I add the figure directly from appendix to here? Or I should always refer to Appendix}.
We then sorted the test set by the highest similarity score for each target figure and selected the top 25\% of data points with high similarity in both metrics, creating the Context-Aligned set with 2,513 target figures.
The remaining 9,690 target figures in the test set formed the Context-Misaligned set.
\autoref{fig:all_figures} shows that personalization is more effective when at least one profile caption is highly similar to the ground-truth caption (Figures~\ref{fig:top_left} and~\ref{fig:bottom_left}).
While profiles still help in the Context-Misaligned set (Figures~\ref{fig:top_right} and~\ref{fig:bottom_right}), their impact is noticeably smaller. 
The detailed results can be found in \autoref{app:easy-hard}.





Graph plot show the best result, also with plenty of sample. Bar chart and scatterplot very similar. Scatterplot is the wrost one(Table \ref{tab:main-result-byType})
% Please add the following required packages to your document preamble:
% \usepackage{multirow}
\begin{table}[t]
\centering
\footnotesize
\renewcommand{\arraystretch}{1.2}
\setlength{\tabcolsep}{2pt}
\begin{tabular}{cccccccc}
\hline
\multirow{2}{*}{\textbf{Figure Type}} & \multicolumn{4}{c}{\textbf{BLEU}} & \multicolumn{3}{c}{\textbf{ROUGE}} \\ \cline{2-8} 
 & \textbf{B-1} & \textbf{B-2} & \textbf{B-3} & \textbf{B-4} & \textbf{R-1} & \textbf{R-2} & \textbf{R-L} \\ \hline
Graph Plot \\ (n=93,900) & .316 & .231 & .184 & .151 & .425 & .235 & .356 \\ \hline
Node Diagram \\ (n=20,520) & .275 & .205 & .164 & .135 & .371 & .205 & .327 \\ \hline
Scatterplot \\ (n=12,024) & .296 & .208 & .161 & .129 & .409 & .210 & .328 \\ \hline
Bar Chart \\ (n=10,285) & .302 & .217 & .170 & .138 & .414 & .219 & .352 \\ \hline
Equation \\ (n=8,772) & \multicolumn{1}{l}{.280} & \multicolumn{1}{l}{.211} & \multicolumn{1}{l}{.170} & \multicolumn{1}{l}{.141} & \multicolumn{1}{l}{.381} & \multicolumn{1}{l}{.214} & \multicolumn{1}{l}{.335} \\ \hline
\end{tabular}
\caption{Performance comparison of LLMs among different figure type as target figure.}
\label{tab:main-result-byType}
\end{table}


Economic seems work best, but the sample is also the least.(Table \ref{tab:main-result-byCategory})
% Please add the following required packages to your document preamble:
% \usepackage{multirow}
\begin{table}[t]
\centering
\footnotesize
\renewcommand{\arraystretch}{1.2}
\setlength{\tabcolsep}{2pt}
\begin{tabular}{cccccccc}
\hline
\multirow{2}{*}{\textbf{Paper Category}} & \multicolumn{4}{c}{\textbf{BLEU}} & \multicolumn{3}{c}{\textbf{ROUGE}} \\ \cline{2-8} 
 & \textbf{B-1} & \textbf{B-2} & \textbf{B-3} & \textbf{B-4} & \textbf{R-1} & \textbf{R-2} & \textbf{R-L} \\ \hline
\begin{tabular}[c]{@{}c@{}}Physics \\ (n=73,452)\end{tabular} & .300 & .212 & .164 & .132 & .414 & .215 & .333 \\ \hline
\begin{tabular}[c]{@{}c@{}}Comp Sci\\ (n=56,988)\end{tabular} & .312 & .236 & .192 & .160 & .413 & .238 & .365 \\ \hline
\begin{tabular}[c]{@{}c@{}}Maths\\ (n=6,120)\end{tabular} & .310 & .236 & .195 & .163 & .405 & .238 & .361 \\ \hline
\begin{tabular}[c]{@{}c@{}}EE and SS\\ (n = 4,092)\end{tabular} & .326 & .249 & .205 & .172 & .425 & .253 & .378 \\ \hline
\begin{tabular}[c]{@{}c@{}}Stats\\ (n=3,228)\end{tabular} & \multicolumn{1}{l}{.302} & \multicolumn{1}{l}{.224} & \multicolumn{1}{l}{.181} & \multicolumn{1}{l}{.151} & \multicolumn{1}{l}{.414} & \multicolumn{1}{l}{.231} & \multicolumn{1}{l}{.349} \\ \hline
\begin{tabular}[c]{@{}c@{}}Quant Bio \\ (n=1,500)\end{tabular} & \multicolumn{1}{l}{.289} & \multicolumn{1}{l}{.208} & \multicolumn{1}{l}{.166} & \multicolumn{1}{l}{.135} & \multicolumn{1}{l}{.403} & \multicolumn{1}{l}{.215} & \multicolumn{1}{l}{.329} \\ \hline
\begin{tabular}[c]{@{}c@{}}Quant Fina \\ (n=900)\end{tabular} & \multicolumn{1}{l}{.333} & \multicolumn{1}{l}{.255} & \multicolumn{1}{l}{.210} & \multicolumn{1}{l}{.176} & \multicolumn{1}{l}{.433} & \multicolumn{1}{l}{.256} & \multicolumn{1}{l}{.376} \\ \hline
Econ (n=169) & \multicolumn{1}{l}{.376} & \multicolumn{1}{l}{.293} & \multicolumn{1}{l}{.244} & \multicolumn{1}{l}{.203} & \multicolumn{1}{l}{.488} & \multicolumn{1}{l}{.302} & \multicolumn{1}{l}{.425} \\ \hline
\end{tabular}
\caption{Performance comparison of LLMs among different paper categories.}
\label{tab:main-result-byCategory}
\end{table}


\paragraph{Using Profile is the best!}

\paragraph{Personalization worked better for the easy data portion.}
\kenneth{TODO Aashish: Please write how you selected the data here. Distribution charts might need to go to Appendix.}
%\kenneth{TODO: (Avg, Max) x (BLEU, BertScore, ROUGE) scores + distribution}

% \tong{to}
\paragraph{Task Difficulty Analysis}
In order to understand how variations in content and vocabulary between the profile captions and target caption affect personalization, we construct 2 subsets: {\dataset} Easy and {\dataset} Hard. For each datapoint present in our test set, we calculate the semantic and lexical similarity of profile captions with the target caption using BERTScore \cite{Zhang*2020-xq} and ROUGE-L \cite{Lin2004-ox} respectively as presented in Figure \ref{fig:profile_target_distribtions}. The test set sorted by the max scores observed between the target caption and any of its profile captions, followed by the sampling of the top 25\% of datapoints exhibiting the highest similarity in each of the 2 metrics. Finally, we select the intersection of datapoints between sets as {\dataset} Easy with the rest constituting {\dataset} Hard. Across both subsets, we observe a consistent increase in caption alignment with the addition of more profiles as presented in Tables \ref{tab:main-result-easy}, \ref{tab:main-result-hard}. However, the improvements found with profile captions on the {\dataset} Hard subset are much more limited. This indicates the challenging nature of this task in utilizing style features to improve figure caption personalization.



Seems 4o the score is similar, while other LLM experience poor result (Table \ref{tab:main-result-hard})




More fair comparison between 1 profile vs all profile if only consider data with 2 or more available profiles (Table \ref{tab:main-result-2more})
% % Please add the following required packages to your document preamble:
% % \usepackage{multirow}
% \begin{table}
% \footnotesize
% \setlength{\tabcolsep}{2pt}
% \centering
% \begin{tabular}{@{}l@{\hspace{1pt}}l@{\hspace{1pt}}cccc@{\hspace{2pt}}ccc@{}}
% \toprule
% \multirow{2}{*}{\textbf{LLM}} & \multirow{2}{*}{\textbf{\begin{tabular}[c]{@{}l@{}}Profile\\ Used\end{tabular}}} & \multicolumn{4}{c}{\textbf{BLEU}} & \multicolumn{3}{c}{\textbf{ROUGE}} \\ \cline{3-9} 
%  &  & \textbf{B-1} & \textbf{B-2} & \textbf{B-3} & \textbf{B-4} & \textbf{R-1} & \textbf{R-2} & \textbf{R-L} \\ \hline
% \multirow{2}{*}{\textbf{\begin{tabular}[c]{@{}l@{}}GPT-\\ 4o\end{tabular}}} & \textbf{One} & .281 & .189 & .139 & .105 & .386 & .181 & .315 \\ \cline{2-9} 
%  & \textbf{All} & .309 & .215 & .163 & .126 & .413 & .209 & .342 \\ \hline
% \multirow{2}{*}{\textbf{\begin{tabular}[c]{@{}l@{}}Llama \\ 4 Scout\end{tabular}}} & \textbf{One} & .373 & .295 & .248 & .212 & .482 & .303 & .425 \\ \cline{2-9} 
%  & \textbf{All} & .419 & .343 & .294 & .258 & .524 & .349 & .468 \\ \hline
% \multirow{2}{*}{\textbf{\begin{tabular}[c]{@{}l@{}}Gemini\\ 2.5 Flash\end{tabular}}} & \textbf{One} & .372 & .293 & .245 & .211 & .483 & .303 & .426 \\ \cline{2-9} 
%  & \textbf{All} & .419 & .342 & .294 & .257 & .526 & .354 & .473 \\ \hline
% \multirow{2}{*}{\textbf{\begin{tabular}[c]{@{}l@{}}GPT-\\ 4o mini\end{tabular}}} & \textbf{One} & .287 & .204 & .157 & .123 & .399 & .209 & .328 \\ \cline{2-9} 
%  & \textbf{All} & .318 & .236 & .188 & .153 & .430 & .245 & .360 \\ \hline
% \end{tabular}
% \caption{Performance comparison of LLMs with different profile settings. On data only with 2 or more profiles available (n=6284)}
% \label{tab:main-result-2more}
% \end{table}


\begin{table}
\centering
\footnotesize
\renewcommand{\arraystretch}{1.2}
\setlength{\tabcolsep}{2pt}
\definecolor{lightgreen}{rgb}{0.9, 1.0, 0.9}
\definecolor{darkgreen}{rgb}{0.7, 0.9, 0.7}
\begin{tabular}{lcccccccc}
\toprule
\multirow{2}{*}{\textbf{LLM}} & \multirow{2}{*}{\textbf{Profile}} & \multicolumn{4}{c}{\textbf{BLEU}} & \multicolumn{3}{c}{\textbf{ROUGE}} \\ 
\cmidrule(lr){3-6} \cmidrule(lr){7-9}
 & \textbf{Used} & \textbf{B-1} & \textbf{B-2} & \textbf{B-3} & \textbf{B-4} & \textbf{R-1} & \textbf{R-2} & \textbf{R-L} \\ 
\midrule
\multirow{2}{*}{GPT-4o} 
 & \cellcolor{lightgreen}One & \cellcolor{lightgreen}.281 & \cellcolor{lightgreen}.189 & \cellcolor{lightgreen}.139 & \cellcolor{lightgreen}.105 & \cellcolor{lightgreen}.386 & \cellcolor{lightgreen}.181 & \cellcolor{lightgreen}.315 \\ 
 & \cellcolor{darkgreen}All & \cellcolor{darkgreen}.309 & \cellcolor{darkgreen}.215 & \cellcolor{darkgreen}.163 & \cellcolor{darkgreen}.126 & \cellcolor{darkgreen}.413 & \cellcolor{darkgreen}.209 & \cellcolor{darkgreen}.342 \\ 
\midrule
\multirow{2}{*}{Llama-4 Scout}
 &  \cellcolor{lightgreen}One & \cellcolor{lightgreen}.373 & \cellcolor{lightgreen}.295 & \cellcolor{lightgreen}.248 & \cellcolor{lightgreen}.212 & \cellcolor{lightgreen}.482 & \cellcolor{lightgreen}.303 & \cellcolor{lightgreen}.425 \\ 
 & \cellcolor{darkgreen}All & \cellcolor{darkgreen}.419 & \cellcolor{darkgreen}.343 & \cellcolor{darkgreen}.294 & \cellcolor{darkgreen}.258 & \cellcolor{darkgreen}.524 & \cellcolor{darkgreen}.349 & \cellcolor{darkgreen}.468 \\ 
\midrule
\multirow{2}{*}{\begin{tabular}[c]{@{}l@{}}Gemini-2.5\\Flash Preview\end{tabular}} 
 & \cellcolor{lightgreen}One & \cellcolor{lightgreen}.372 & \cellcolor{lightgreen}.293 & \cellcolor{lightgreen}.245 & \cellcolor{lightgreen}.211 & \cellcolor{lightgreen}.483 & \cellcolor{lightgreen}.303 & \cellcolor{lightgreen}.426 \\ 
 & \cellcolor{darkgreen}All & \cellcolor{darkgreen}.419 & \cellcolor{darkgreen}.342 & \cellcolor{darkgreen}.294 & \cellcolor{darkgreen}.257 & \cellcolor{darkgreen}.526 & \cellcolor{darkgreen}.354 & \cellcolor{darkgreen}.473 \\ 
\midrule
\multirow{2}{*}{GPT-4.1 Mini}  
 & \cellcolor{lightgreen}One & \cellcolor{lightgreen}.287 & \cellcolor{lightgreen}.204 & \cellcolor{lightgreen}.157 & \cellcolor{lightgreen}.123 & \cellcolor{lightgreen}.399 & \cellcolor{lightgreen}.209 & \cellcolor{lightgreen}.328 \\ 
 & \cellcolor{darkgreen}All & \cellcolor{darkgreen}.318 & \cellcolor{darkgreen}.236 & \cellcolor{darkgreen}.188 & \cellcolor{darkgreen}.153 & \cellcolor{darkgreen}.430 & \cellcolor{darkgreen}.245 & \cellcolor{darkgreen}.360 \\ 
\bottomrule
\end{tabular}
\caption{Performance comparison of LLMs with different profile settings. On data only with 2 or more profiles available (n=6284)}
\label{tab:main-result-2more}
\end{table}


Under same 1 profile configuration, when target figure type is same as profile figure type, the performance is better (Table \ref{tab:main-result-sameType})



---standard reference-based metrics that measure how closely the generated captions match the original author-written captions. 
The results show that incorporating profile information consistently improves caption quality across all four models. Additionally, using all profile figures provides better results than using just one.




\kenneth{--------------------KENNETH IS WORKING HERE--------------------}

%\kenneth{probably need to use a smaller model}
For our caption generation task, we selected four models with varying capabilities: (1) GPT-4o, (2) Llama 4 Scout, (3) Gemini 2.5 Flash Preview, and (4) GPT-4.1 Mini. \ed{add model citation.} \ed{ Specifically, we utilize OpenRouter\footnote{\url{https://openrouter.ai}}, a unified API that grants access to a wide range of AI models via a single endpoint, while automatically managing fallback mechanisms and optimizing for cost-effectiveness}. This diverse selection allowed us to test whether our proposed profile configuration produces consistent improvements across different model architectures, thereby validating that the personalized approach enhances scientific figure caption generation regardless of the underlying model used.


\paragraph{Experiment Setups.}
\kenneth{explain why use automatic and why we didn't normalized}

For each model, we have 3 different configurations to generate captions: (1) No Profile, (2) One Profile, (3) All Profile. 
%The detailed of the input difference can checked the Table~\ref{tab:config}
%\lee{what are inputs exactly? Unclear to me. Do you use a target figure and profile figures (and their captions) to generate the caption of the target figure? If so, spell out what are inputs and outputs}
%\sam{I just add the table to illustrate, see if helps} \ed{add some explanations for the three settings}

\begin{table}[t]
\centering
\small
\begin{tabular}{lll}
\toprule
\textbf{Config.} & \textbf{Target} & \textbf{Profile} \\ \hline
\multicolumn{1}{l|}{No-Profile} & \multicolumn{1}{l|}{Img + Para} & None \\ \hline
\multicolumn{1}{l|}{One-Profile} & \multicolumn{1}{l|}{Img + Para} & 1 profile (Img + Para + Cap) \\ \hline
\multicolumn{1}{l|}{All-Profile} & \multicolumn{1}{l|}{Img + Para} & All profiles (Img + Para + Cap) \\ 
\bottomrule
\end{tabular}
\caption{Input settings for each configuration in the generation task.}
\label{tab:config}
\end{table}
For the data output cleaning, here are the steps:
1. Use GPT-4o-mini to clean the output (e.g. clean the reasoning steps and irrelevant pre- and suf-fix of the "real context" caption from the response of GPT such as figure index included in the output) --> may include the prompts used to clean the output 

2. First round of cleaning we found that there was 11 case (out of 12,270) with the issue of the original caption as there was some cases such as parsing error such as capturing the figure caption without real context (just figure index) or capturing the wrong caption, so the comparison between the generated output with these case would be invalid 

3. Second round, we filtered out the response that the generation was failed (Appedix table?), including the blank cases. Unique case clean across all model and configuration: 56 (out of 12,259)

For the evaluation, we adopted automatic evaluation by calculating BLEU and ROUGE score to compare the similarity of generated caption to the original captions of the target figures.
    
\end{comment}
